\documentclass[border=4pt]{standalone}
\usepackage{amsmath,amssymb,mathtools,xcolor,varwidth}
\definecolor{tanorange}{HTML}{E8770A}
\definecolor{arcblue}{HTML}{1E6FE0}
\definecolor{chordgreen}{HTML}{00A35A}
\definecolor{anglered}{HTML}{D81B60}
\begin{document}
\begin{varwidth}{28em}
\large\bfseries
Draw \textcolor{tanorange}{the tangent $AD$}, the straight line that just grazes \textcolor{arcblue}{the curve} at $A$, matching its direction there and extending both ways. Between \textcolor{tanorange}{the tangent} and \textcolor{chordgreen}{the chord} opens \textcolor{anglered}{the angle $BAD$}. As $B$ approaches $A$, \textcolor{chordgreen}{the chord} swivels toward \textcolor{tanorange}{the tangent}, and \textcolor{anglered}{this angle} is diminished \emph{in infinitum}. For if \textcolor{anglered}{the angle} did \emph{not} vanish, \textcolor{arcblue}{the arc} would meet \textcolor{tanorange}{the tangent} at a finite rectilinear angle — a sharp elbow at $A$ — and the curvature there would not be continuous, against the supposition. Hence ultimately $\angle BAD \to 0$. \textit{Q.E.D.}
\end{varwidth}
\end{document}
